\chapter{Introduction}

\section{What the heck is a plasma}

\subsection{Plasmas in the universe}

\subsection{How do plasmas affect us on earth?}

\subsection{Where does magnetic reconnection play into this whole shabang?}

\section{Magnetic Reconnection}

\subsection{The Generalized Ohm's Law}

\subsubsection{How do we derive this?}

\subsubsection{A breakdown of the terms}

\subsection{Regimes of Reconnection}

\subsubsection{Reconnection scaling parameters}

i.e. why do we use $S$ and $L / d_i$ to characterize reconnection regimes?

\subsubsection{Where do various systems line up?}

\subsubsection{What are we interested in?}

\begin{figure}
  \caption{Diagram of reconnection regimes.}
\end{figure}

\subsection{Sweet-Parker Reconnection}

\subsubsection{A quick historical detour}

\subsubsection{What are the assumptions that go into Sweet-Parker?}

\subsubsection{Sweet-Parker Reconnection Rate Derivation}

\begin{figure}
  \caption{Diagram of Sweet-Parker reconnection.}
\end{figure}

\subsubsection{Application to Real Systems}

\paragraph{What are some example parameters where reconnection occurs?}

\paragraph{What time scales does Sweet-Parker predict?}

\paragraph{What did we miss?}

\subsection{Hall Reconnection}

\subsubsection{Why should the Hall term appear?}

\paragraph{What is the relevant physics that sets when the hall term matters}

\paragraph{In what kind of plasmas does the hall term appear and when does it not?}

\begin{figure}
  \caption{Diagram of 2d hall reconnection}
\end{figure}

\subsection{Kinetic Reconnection}

\subsubsection{What makes kinetic reconnection ``kinetic''?}

\subsubsection{Where does Pressure Anisotropy come from?}

\paragraph{Movement of flux tube}

\begin{figure}
  \caption{Flux tube changing size during reconnection causing pressure anisotropy build up.}
\end{figure}

\paragraph{Effect of background field}

\subparagraph{Too low of background field does...}

\subparagraph{Too high of field does\dots}

\subparagraph{Goldilocks!}

\begin{figure}
  \caption{Simulated kinetic reconnection with changing guide field.}
\end{figure}

\subsubsection{How does pressure anisotropy change what the reconnection looks like}

\begin{figure}
  \caption{Diagram from Ohia of with and without pressure anisotropy in simulation.}
\end{figure}

\paragraph{Thinner reconnection layer?}

Use results from Jan's more recent papers.

\begin{figure}
  \caption{Some figure from Jan's papers thowing effects of pressure anisotropy}
\end{figure}

\paragraph{Marginal Firehose Condition}

\paragraph{Concluding: what features should we look for in the data?}

\subsubsection{Why do we want to measure it?}

\section{Experimental Setup}

\subsection{Overall how does our experiment progress}

\subsection{Initial field}

\subsubsection{Helmholtz}

\subsubsection{Permanent magnets}

\subsubsection{Toroidal}

\paragraph{Shifted coil}

\paragraph{Axial coil}

\subsubsection{Perturbation coils}

\subsection{Plasma Generation}

\subsubsection{How do our guns work}

\paragraph{Gas lines}

\paragraph{PFNs}

\begin{figure}
  \caption{Circuit diagram of plasma gun.}
\end{figure}

\paragraph{Puff trigger}

\paragraph{How to breakdown}

\begin{figure}
  \caption{Simplified cut of plasma gun insides.}
\end{figure}

\subsubsection{How are the guns placed for injection?}

\subsubsection{What goes wrong with them}

\paragraph{Burning up}

\begin{figure}
  \caption{Image from high speed camera of gun spewing}
\end{figure}

\paragraph{Burn holes in wall because of background field}

Get image from Paul with hole in wall.

\paragraph{Air cooling}

\paragraph{Optimizing gun operation}

\subsection{Drive field}

\subsubsection{Individual coils}

\paragraph{Zone's between coils as plasma sources}

\subsubsection{Drive Cylinder}

Reference Paul's RSI + thesis.

\subsection{Diagnostics Overview}

\begin{table}
  \caption{List of diagnostics used in the experiment and what they can measure.}
\end{table}

\subsection{How does this compare to space}

\begin{table}
  \caption{Comparison of lab and space parameters.}
\end{table}
